% Français 9115, batch 31 — Gallica views 601–620.
% Pass: Opus 5.5 (claude-opus-5-5), 2026-09-22 — first pass, unchecked against the views by a human.
%
% What the twenty views hold. Two pieces, both hers as far as the leaves
% show, both on vibrating elastic surfaces; none of it number theory.
%
% (1) v. 601–613: eleven pages of the paginated memoir on vibrating
% plates, in her regular fair hand, lightly corrected, pages (62)–(72) at
% the head, centre, written on both sides of each leaf, the ink series on
% the side whose page number is odd: 296 (v. 602, (63)), 297 (604, (65)), 298 (606, (67)), 299 (608,
% (69)), 300 (610, (71)). It is the comparison of theory with Chladni's
% experiments on square plates — the same section, often sentence for
% sentence, that batch 25 read at v. 488–500 under pages (78)–(90): the
% same N.o 16 « comparaison des sons », the same references to N.o 13 and
% N.o 15, the same tables. So the volume holds two redactions of it, some seventeen
% pages apart in her numbering; which is the earlier the leaves here do
% not say. Beginning and end are both outside the batch.
%   v. 601  (62) integral (I) with b′ = 0, n = 1, m = 5: the factor sin(πx/A)
%           gives a nodal line through the middle, the other factor two
%           fixed points of rest; Chladni's fig. 79 b fits except that three
%           lower lines are replaced by an « equivalent » figure; then
%           m = 6 (fig. 86), the formula left unfinished at the edge.
%   v. 602  (63) fig. 86 and its two re-entrant curves, taken for an
%           accident of the plate; her application of (I) to figs 79 b and
%           86 « n'est que conjecturale »; sounds ∝ 25 and 36; more than
%           half of Chladni's square-plate figures accounted for; she will
%           follow Chladni's equivalences for the rest.
%   v. 603  (64) N.o 16 « comparaison des sons »: Chladni's own bound on
%           experimental error (« p. 154 de son livre », quoted: never more
%           than a semitone), hence at most a tone between two ratios; the
%           « vibrations tournantes » first.
%   v. 604  (65) the theoretical numbers for figs 63, 66 a, 69–70, 74 a,
%           79 a (2, 5, 10, 17, 26), their ratios to fig. 63, and Chladni's
%           « table p. 152 du traité d'acoustique »; 63 to 66 a and 63 to
%           69 agree; fig. 70 a tone off fig. 69, unexplained.
%   v. 605  (66) 63 to 74 within the error; 63 to 79 a « beaucoup trop
%           petit », about a tone too sharp; 79 b worse; 63 to 86, 64/3
%           against 18.
%   v. 606  (67) figs 79 b and 86 only « hypothétique »; the figures with
%           complete periods in both directions, compared with fig. 68 a
%           (8): figs 71, 75, 80 a, 81 a, 82, 87–88 (13, 18, 25, 29, 32, 40),
%           ratios, Chladni's sounds.
%   v. 607  (68) each ratio checked: 71, 75, 82, 87–88 agree, 80 by a
%           comma, 81 by more than a comma; 68 to 79 within the comma
%           128/125.
%   v. 608  (69) Chladni's notation m|n for the cases he gave no figure
%           for; theoretical numbers for 2|2 … 6|5 (8, 20, 34, 45, 41, 52,
%           50, 61), their ratios, Chladni's sounds; 4|2 agrees.
%   v. 609  (70) 5|3, 6|3, 5|4, 6|4, 5|5, 6|5 each within a semitone or a
%           comma; « Il est impossible d'espérer plus d'accord… ».
%   v. 610  (71) the same agreement for the cases Chladni's authority
%           assigned; « Cependant il se présente ici une difficulté »:
%           against the sounds of integral (C) and of the free lamina
%           (d²z/dx² = 0, d³z/dx³ = 0), every figure with complete periods
%           is about a tone too grave — 0|4 (figs 72–73) sol♯ 3 against
%           4|3 (fig. 80) fa♯ 5, where Euler's theory and hers want two
%           octaves (25/4 and 25).
%   v. 613  (72) that tone is a minimum, approached the more closely by the
%           « vibrations tournantes » the more nodal lines they have (figs
%           63, 68, 66 a, 69, 74, 79 a); the same on rectangular plates.
%           The page stops mid-sentence (« dans la comparaison au son des
%           figures »); its sequel is not in this batch.
%   Three passages, those on figs 79 b and 86 (v. 601, 602, 605–606) and
%   on 68–79 (v. 607), are crossed by a single long thin oblique stroke,
%   and the matching lines of the tables (v. 604, 606) are struck: she
%   seems to have withdrawn figs 79 b and 86, and 79 a from the 68 a
%   comparison. The text is transcribed as written, with a note each time.
% (2) v. 614–620: four leaves of a different paper, unpaginated, in a
%   faster hand heavily corrected between the lines, inked 301–304 at the
%   top right: drafts, it seems, of the opening of a prize memoir on
%   elastic surfaces addressed to « la classe ».
%   v. 614  (301) from mid-sentence: the principal curvatures as
%           coordinates; the equation of vibrating elastic surfaces is of the
%           fourth degree, can be lowered to the second, and one of its
%           factors is the equation of tense membranes, so that their
%           particular integrals serve the elastic surface too.
%   v. 616  (302) « Les geometres ne se sont pas encore beaucoup occupés »
%           of these integrals; « M.r Biot est je crois le seul » to have
%           shown nodal lines parallel to the sides of rectangular
%           membranes; she will show it follows from the periodicity; the
%           « singularité remarquable »: only where elastic and tense
%           surfaces share a figure may the vibrating sheets be likened to
%           Sauveur's waves and nodes.
%   v. 618  (303) continues 616: that analogy, « indiquée par les termes
%           mêmes du programme publié par la classe », fails elsewhere; the
%           plan (the differential equation, then particular integrals
%           compared with experiment); « je m'estimerai heureux si la classe
%           juge que j'ai reussi a ebaucher une théorie que des » — and
%           stops.
%   v. 620  (304) the cylindrical surface on a circular base, equation (e)
%           with the term (1/R²)(d²z″/dx² + d²z″/dy²); the portion of arc
%           A, R = A/a; such surfaces « n'a encore, je crois, été l'objet
%           d'aucune experience »; struck: « Le tems m'a manqué… »; then
%           « Je me suis résolue beaucoup trop tard a reprendre le sujet »,
%           time lacking both for the experiments and for the integrals;
%           she will first give the general results of the experiments.
%           Stops on « qu'ils présente ».
%   The adjectives of the author shift: « parvenu », « heureux » (v. 618)
%   are masculine, « privée », « résolue » (v. 620) feminine.
%
% Views not transcribed. v. 611: the page of v. 609 photographed a second
% time (same « (70) », same text, framing differs). v. 612: the page of
% v. 610 a second time (same « (71) », same ink 300, same text). v. 615,
% 617, 619: the blank backs of the leaves of v. 614, 616, 618, the recto
% showing through reversed (checked at a raised contrast). v. 621 (not
% transcribed, glanced at) is the blank back of v. 620.
%
% Seams. At 600/601 nothing is cut: view 600 (glanced at, its head not read) ends
% on « m = 5, m = 6 (en prenant pour … des valeurs convenables) » and
% view 601 opens « En effet, si on fait b′ = 0, n = 1 et m = 5 ». At
% 620/621 the draft stops mid-sentence on « des singularités qu'ils
% présente »; v. 621 is blank, so any sequel is in batch 32. The fair
% memoir stops at v. 613, page (72), mid-sentence; its page (73) is not in
% this batch.
%
% The numbers on the leaves, and why no \folio{} is written. (1) Her
% pagination, in parentheses at the head, centre: (62) v. 601, (63) 602,
% (64) 603, (65) 604, (66) 605, (67) 606, (68) 607, (69) 608, (70) 609 and
% its repeat 611, (71) 610 and its repeat 612, (72) 613 — without a gap;
% (63), (65), (67), (69), (71) — the pages that carry the ink number — are
% struck through with a horizontal stroke, the others not. (2) The ink
% series at the top right, in the larger rounder hand: 296 (v. 602), 297
% (604), 298 (606), 299 (608), 300 (610, and 612), 301 (614), 302 (616),
% 303 (618), 304 (620) — one per leaf, never on a back; it continues 295 at v. 600,
% as batch 30 reads it (that file appeared after this pass had read its
% own views; its page (61) at v. 600 fits this batch's (62) at v. 601). Penned, as everywhere in this volume, so no \folio{}. (3) On
% v. 620, beside 304, a small blotted figure, perhaps 1, struck and
% underlined. (4) Her article number N.o 16 (v. 603), and references to
% N.o 13, N.o 15 (v. 604, 606); integrals (C), (I), equation (e). (5)
% Chladni's figure numbers (63–88 with a, b) and his pages « p. 152 »,
% « p. 154 ». Every number above was read on its view; none is inferred.
% None of these leaves can be Laubenbacher and Pengelley's « Manuscript C »
% (ff. 348r–349r): the subject is elastic plates, and the ink series here
% runs 296–304.
%
% Notation. Hers, kept. A the side of the square plate; m, n, b′ the
% parameters of her integrals; her cursive capital beside a′ is set
% $\mathcal{A}$, as batch 22 did. Chladni's m|n (m straight nodal lines
% parallel to one axis, n to the other) keeps her vertical stroke. Notes by
% name with the octave number after them, the sharp $\sharp$ and flat
% $\flat$ written above the note (« si♭ 3 + grave »). Her braces are kept.
% « sin », « cos » set \text{sin}, \text{cos}. Chladni's lists and her
% tables of ratios are set as arrays, with her dotted leaders. Words she
% underlines in prose are set in \emph{}; letters underlined inside
% mathematics keep \underline{}. Where she cancels a plural with her small
% cross (« lex sonx », v. 613; « valeurs », v. 614) the singular is given,
% with a note. Spelling is hers and is not flagged: « complettes »,
% « parallelles », « précédans », « suivans », « quarré », « tems »,
% « dela », « delaplaque », « desorte », « ensorte », « aulieu »,
% « audessous », « cequi », « parceque », « c'est adire », « leplus »,
% « emploirai », « signifira », « continurai », « Aucontraire »,
% « inapréciable », « bizares », « crû », « pû », « dueau »,
% « l'expériencedans », « ladifficulté ».
%
% What could not be read, and what is offered with doubt. \ill{} stands
% five times: the struck first word of v. 614; a struck word after « des
% mêmes cas » (v. 614); in v. 620 the small sign under the exponent of b
% in both equations, and the blotted sign before 1/R² in the first.
% \uncertain{} stands for « 69 » in the list of v. 604 (perhaps 60),
% « courbures », « valevoir » (va le voir), « delà » (v. 614), three
% struck groups of v. 616 (« ni en general », « ce savant », « d'après
% les le plus »), « résolue » (v. 620). Nothing was guessed to fill a gap.
%
% Nothing here was checked against any published edition of Sophie
% Germain's papers, of Chladni, or of any memoir on plates.
\documentclass[11pt,a4paper]{article}
% The preamble shared by every transcription of the mathematicians' manuscripts in Gallica.
%
% It rests on one idea: **uncertainty compounds, it does not resolve.** A
% transcription that smooths over a doubtful word has destroyed the only thing
% it was made for — a reader must be able to tell, at every point, what was on
% the page and what was supplied. The apparatus macros below are therefore the
% critical apparatus, not decoration.
%
% The same commands are recognised by `scripts/render.mjs`, which produces the
% site's reading view, and by `scripts/tei.mjs`. A macro added here must be
% added there too, or rendering fails loudly — which is the wanted behaviour.
%
% Comments here are English, like the rest of the repository. The documents
% this preamble sets are French, and that is deliberate: see the skills.

\ProvidesPackage{germain}[2026/09/15 Transcriptions of mathematicians' manuscripts in Gallica]

\usepackage{amsmath,amssymb,amsthm}
\usepackage{mathrsfs}
% Kept from the parent preamble so the renderer's diagram support stays
% available; these manuscripts draw no commutative diagrams, and nothing here uses it.
\usepackage{tikz-cd}
\usepackage[english,french]{babel}
\usepackage{fontspec}

% --- Greek ------------------------------------------------------------------
% The text face stays fontspec's default, Latin Modern, which has no Greek:
% XeTeX drops every glyph it lacks with a line in the log and nothing on the
% page, and the Greek words the manuscripts quote vanished from the PDFs. So
% the Greek and Greek Extended blocks get a XeTeX character class of their own,
% and entering or leaving a run of it switches the family to CMU Serif — the
% same Computer Modern design, with polytonic Greek — and back. Only the
% family changes: a Greek word in \textit or \textbf keeps its shape and series.
% The transitions to and from every other class carry the tokens that class
% has towards an ordinary letter, so French spacing before « ; » or inside
% « » still holds after a Greek word. No grouping, so no brace or \endgroup
% can come out unbalanced; maths is left alone, since XeTeX inserts nothing
% there. Kept identical in `exercices.sty`.
\makeatletter
\newfontfamily\germainGreekfont{cmun}[
  Extension=.otf, NFSSFamily=germaingreek,
  UprightFont=*rm, ItalicFont=*ti, SlantedFont=*sl,
  BoldFont=*bx, BoldItalicFont=*bi, BoldSlantedFont=*bl]
\newXeTeXintercharclass\germain@greek
\def\germain@greekrange#1#2{%
  \@tempcnta=#1\relax
  \loop
    \XeTeXcharclass\@tempcnta=\germain@greek
    \ifnum\@tempcnta<#2\advance\@tempcnta\@ne\repeat}
\germain@greekrange{"0370}{"03FF}
\germain@greekrange{"1F00}{"1FFF}
\def\germain@greekprev{\rmdefault}
\def\germain@greekin{%
  \edef\germain@greekprev{\f@family}\fontfamily{germaingreek}\selectfont}
\def\germain@greekout{\fontfamily{\germain@greekprev}\selectfont}
\def\germain@greekjoin{%
  \XeTeXinterchartokenstate=\@ne
  \@tempcnta=\z@
  \loop
    \ifnum\@tempcnta=\germain@greek\else
      \XeTeXinterchartoks\germain@greek\@tempcnta=\expandafter{%
        \expandafter\germain@greekout\the\XeTeXinterchartoks\z@\@tempcnta}%
      \XeTeXinterchartoks\@tempcnta\germain@greek=\expandafter{%
        \the\XeTeXinterchartoks\@tempcnta\z@\germain@greekin}%
    \fi
    \ifnum\@tempcnta<4095\advance\@tempcnta\@ne\repeat}
% babel-french sets its own classes, and saves and restores the interchar
% state, each time a language is selected: join again after it does.
\addto\extrasfrench{\germain@greekjoin}
\addto\extrasenglish{\germain@greekjoin}
\AtBeginDocument{\germain@greekjoin}
\makeatother

\usepackage{geometry}
\geometry{a4paper,margin=25mm,marginparwidth=28mm}
\usepackage{marginnote}
\usepackage{xcolor}
\usepackage[normalem]{ulem}
\setlength{\emergencystretch}{3em}
\usepackage{hyperref}
\hypersetup{colorlinks=true,linkcolor=black,urlcolor=black,pdfborder={0 0 0}}

% --- Batch metadata ---------------------------------------------------------
% Taken as they stand by the rendering script, which shows them at the head of
% the reading view. They fix what the file claims to cover. The macro names are
% the parent project's, so that the scripts stay shared; \folder here means the
% volume's slug — `fr-9115` — and \pages counts Gallica views.

\newcommand{\folder}[1]{\gdef\grothFolder{#1}}
\newcommand{\batch}[1]{\gdef\grothBatch{#1}}
\newcommand{\pages}[2]{\gdef\grothFirst{#1}\gdef\grothLast{#2}}
\newcommand{\foldertitle}[1]{\gdef\grothFoldertitle{#1}}

% The holder's own shelfmark, printed as they write it: « Français 9115 ».
\newcommand{\shelfmark}[1]{\gdef\grothShelfmark{#1}}

% The dating, copied verbatim from the holder's catalogue — never inferred
% here. « XIXe siècle » is the BnF's; a pass that dates a leaf from its content
% says so in a \note{}, as its own inference.
\newcommand{\dating}[1]{\gdef\grothDating{#1}}

% The status notice, printed in the title block and repeated as a diagonal
% watermark on every page of the PDF. The manuscripts are in the public
% domain; what this line declares is not a right but a quality — an
% unverified first pass by a machine. The skills write the line; a file
% without it gets no watermark, so leaving it out is visible in the output.
\newcommand{\watermark}[1]{\gdef\grothWatermark{#1}}

\gdef\grothFolder{?}\gdef\grothBatch{}\gdef\grothFirst{?}\gdef\grothLast{?}%
\gdef\grothFoldertitle{}\gdef\grothShelfmark{}\gdef\grothDating{}\gdef\grothWatermark{}

% --- The résumé -------------------------------------------------------------
\newenvironment{resume}
  {\small\setlength{\parskip}{0.45em}}
  {\par\medskip\hrule\medskip}

% --- Keywords ---------------------------------------------------------------
% In English, closing the modernised reading's résumé: the modern vocabulary
% under which to search for what the leaves construct. The manifest extracts
% them to tag the volume — this line is the single source of the tags.
\newcommand{\keywords}[1]{\par\smallskip\noindent
  {\footnotesize\textsc{Keywords} --- \textit{#1}}\par}

% --- The critical apparatus -------------------------------------------------

% The Gallica view number, in the margin. This is the anchor that lets a
% disagreement over a formula be settled by looking at *one* image rather than
% a volume of 750. The site uses it to turn the facsimile as the transcription
% is scrolled. A view is one scanned image: usually one side of a leaf.
\newcommand{\page}[1]{%
  \marginnote{\footnotesize\color{gray!70}\textsf{v.\,#1}}%
  \label{page:#1}\ignorespaces}

% The BnF's pencilled foliation, where the leaf carries it and a pass can read
% it: « 198r ». This is what the literature cites — « ff. 198r–208v » — and
% the one line that turns a citation into a view. Recorded, never inferred:
% a folio a pass cannot read is not written.
\newcommand{\folio}[1]{%
  \marginnote{\footnotesize\color{gray!50}\textsf{f.\,#1}}\ignorespaces}

% The modernised reading's coarser counterpart to \page{}: which views a
% section is drawn from.
\newcommand{\pagerange}[2]{%
  \marginnote{\footnotesize\color{gray!70}\textsf{v.\,#1--#2}}%
  \ignorespaces}

% Illegible: never guessed.
\newcommand{\ill}{\textcolor{red!60!black}{[\,\ldots\,]}}

% A doubtful reading: the word is given, and flagged as uncertain.
\newcommand{\uncertain}[1]{\uline{#1}}

% An editorial addition — a missing letter, an implied word.
\newcommand{\add}[1]{\textcolor{blue!50!black}{[#1]}}

% Struck out by the writer. What was crossed out is part of the document.
\newcommand{\struck}[1]{\sout{#1}}

% The transcriber's note, on the state of the page or the mathematics.
\newcommand{\note}[1]{\par\smallskip{\small\color{gray!60!black}\textit{[#1]}}\par\smallskip}

% A marginal note of hers, distinct from the body of the text.
\newcommand{\marginal}[1]{\par\smallskip\noindent\hspace{1em}%
  {\small\color{gray!60!black}#1}\par\smallskip}

% --- Title ------------------------------------------------------------------
% The title block is French, like the documents themselves.

\AddToHook{shipout/background}{%
  \ifx\grothWatermark\empty\else
    \begin{tikzpicture}[remember picture, overlay]
      \node[rotate=52, scale=3.4, align=center, text=gray!18,
            font=\sffamily\bfseries]
        at (current page.center) {\grothWatermark};
    \end{tikzpicture}%
  \fi}

\AtBeginDocument{%
  \begin{center}
    {\large\bfseries Manuscrits de mathématiciens (Gallica) ---
      \ifx\grothShelfmark\empty\grothFolder\else\grothShelfmark\fi}\\[2pt]
    {\itshape\grothFoldertitle}\\[4pt]
    {\small vues \grothFirst--\grothLast
      \ifx\grothBatch\empty\else\ (lot \grothBatch)\fi}\\[2pt]
    \ifx\grothDating\empty\else
      {\small\color{gray!60!black}Datation du catalogue : \grothDating}\\[2pt]
    \fi
    \ifx\grothWatermark\empty\else
      {\small\bfseries\color{red!45!gray}\grothWatermark}\\[2pt]
    \fi
    {\footnotesize\color{gray!60!black}Transcription établie d'après les images IIIF de Gallica.
     Source gallica.bnf.fr / Bibliothèque nationale de France.\\
     Les lectures incertaines sont soulignées, les illisibles notées [\,\ldots\,],
     les ajouts entre crochets ; \textsf{v.} numérote les vues de Gallica,
     \textsf{f.} la foliotation au crayon de la BnF.}
  \end{center}
  \vspace{1em}}

\endinput


\folder{fr-9115}
\shelfmark{Français 9115}
\batch{31}
\pages{601}{620}
\dating{1801-1900}
\watermark{Lecture automatique\\non vérifiée}
\foldertitle{Papiers de Sophie GERMAIN. — Recueil de dissertations et problèmes mathématiques et physiques.}

\begin{document}

\note{Numérotation des feuillets. Les vues 601 à 613 portent en tête, au
milieu, entre parenthèses, la pagination de son mémoire : (62) à la vue
601, puis (63) à (72) aux vues 602 à 613, sans lacune ; les vues 611 et
612 reprennent les pages (70) et (71). Les pages (63), (65), (67), (69) et
(71) sont barrées d'un trait horizontal. Les feuillets sont écrits des
deux côtés, et seul le côté qui porte la page barrée a en haut à droite, à
l'encre, d'une écriture plus ronde, un nombre de la série qui court dans
tout le volume : 296 (vue 602), 297 (604), 298 (606), 299 (608), 300 (610
et 612). Les feuillets de brouillon des vues 614 à 620, sans pagination,
portent 301, 302, 303, 304. Cette série occupe la place de la foliotation
de la Bibliothèque, mais elle est tracée à la plume et non au crayon, et
aucun « f. » n'est donc porté. Tous ces nombres ont été lus sur la vue ;
aucun n'est déduit. Les vues 611 et 612 sont une seconde prise des vues
609 et 610 ; les vues 615, 617 et 619 sont les dos blancs des feuillets
qui les précèdent ; elles ne sont pas transcrites.}

\page{601}

\note{En tête, au milieu, entre parenthèses : « (62) », sans rature. Aucun
nombre en haut à droite : c'est le dos du feuillet de la vue 600, dont le
numéro n'a pas été lu ici. Le bord droit du feuillet passe sous l'onglet de montage et cache la
fin de plusieurs lignes ; les lettres complétées ici entre crochets
d'addition sont celles que le sens rend certaines.}

En effet, si on fait $b'=0$, $n=1$ et $m=5$, l'intégrale (I) deviendra
\[
z=\ldots\,\text{sin}\left(\frac{\pi.x}{A}\right)\Bigl\{a'\,\text{cos}\left(\frac{\pi.x}{A}\right)\Bigl\{1-\frac{10}{3}\,\text{cos}^2\left(\frac{\pi.x}{A}\right)+\frac{4}{3}\,\text{cos}^4\left(\frac{\pi.x}{A}\right)+\&c\Bigr\}+\mathcal{A}\,\text{sin}\left(\frac{\pi.5y}{A}\right)\Bigr\}\,;
\]
\note{Le 4 du second coefficient a la forme de « b » qui est la sienne ; le
même coefficient est écrit $\frac{14}{3}$ quatre lignes plus bas. Le
dénominateur 3 pourrait aussi être lu 9 : ses 3 et ses 9 se ressemblent.}
le facteur $\text{sin}\left(\frac{\pi x}{A}\right)$ donnera dabord une ligne nodale parallele à l'axe
des $y$ ; et cette ligne passera par le milieu delaplaque dont les
extremités sont libres. A ses deux limites paralleles a la ligne nodales,
on aura $\text{cos}\,\frac{\pi x}{A}=0$, desorte que ce facteur fournira à lui seul
deux \emph{points de repos invariables} donnés, comme je l'ai expliqué
plusieurs fois, par l'intersection des lignes nodales qui composent
les figures dues aux suppositions successives $a'=0$ et $\mathcal{A}=0$. Le
facteur $1-\frac{10}{3}\,\text{cos}^2\left(\frac{\pi x}{A}\right)+\frac{14}{3}\,\text{cos}^4\left(\frac{\pi x}{A}\right)+\&c.$ donnera aussi 2 \emph{points de rep\add{os}
invariables}, puisqu'il y a deux valeurs de $\text{cos}\left(\frac{\pi x}{A}\right)$ qui peuvent le ren\add{dre}
nul.
\note{Dans ce second facteur, le 10 est surchargé et son dénominateur
(lu 3) est taché ; l'exposant 2 du premier cosinus paraît écrit sur un
autre chiffre ; le 1 du second coefficient, $\frac{14}{3}$, est bien écrit,
alors que la formule ci-dessus porte $\frac{4}{3}$. Le signe moins
devant $\frac{10}{3}$ est un trait épais.}
Cela posé, il devra y avoir deux flexion complettes su\add{r}
les lignes sinueuses qui composeront la figure. La fig. 79 \emph{b} est
en effet conforme à ces conditions : mais tandis qu'elle devroit mon\add{trer}
5 lignes sinueuses, si la relation entre $a'$ et $\mathcal{A}$ étoit la mê\add{me}
pour les deux moitiés delaplaque, il se trouve que les trois
lignes inferieures sont remplacés par une autre figure nodale q\add{ui}
je crois équivalente.

Si en continuant de faire $b'=0$, $n'=1$ on prend $m=6$, l'équation
(I) deviendra
\[
\begin{array}{l}
z=\ldots\,\text{sin}\left(\frac{\pi.x}{A}\right)\Bigl\{a'\,\text{cos.}\left(\frac{\pi x}{A}\right)\Bigl\{1-\frac{31}{2.3}\,\text{cos}^2\left(\frac{\pi x}{A}\right)+\frac{31.19}{2.3.4.5}\,\text{cos}^4\left(\frac{\pi x}{A}\right)\\
\qquad+\frac{91.19}{2.3.4.5.6.7}\,\text{cos}^6\left(\frac{\pi x}{A}\right)+\&c.+\mathcal{A}\,\text{sin}
\end{array}
\]
\note{La formule est sur une seule ligne sur le feuillet ; elle s'arrête au bord droit, sur « $+\mathcal{A}$ sin »,
et ne reprend pas à la ligne suivante. Les premiers chiffres des trois
numérateurs sont lus 3, 3 et 9 tels qu'ils paraissent ; ses 3 et ses 9 se
confondent, et il se peut que les trois nombres soient le même.}
on trouvera, comme dans le cas précédent, qu'il doit y avoir une lign\add{e}
nodale dueau facteur $\text{sin}\left(\frac{\pi x}{A}\right)$, et que la fonction de $\underline{x}$ que ce fact\add{eur}
multiplie, doit donner quatre \emph{points de repos invariables} ; ensorte
que les lignes sinueuses dont la fig. pourra être composée, devront
également présenter 2 flexion complettes. Mais le terme $\text{sin}\left(\frac{\pi.6y}{A}\right)$ ind
\note{Au premier cas elle écrit « $n=1$ », au second « $n'=1$ ». « dueau » : « dû au », tel qu'écrit. Les mots soulignés par elle sont
donnés en italique. La page s'arrête sur « ind », coupé par l'onglet ;
la suite est au haut de la vue 602. Un long trait fin, presque vertical,
descend à travers le texte depuis la ligne « A ses deux limites » jusqu'à
« La fig. 79 \emph{b} est » ; il ne suit aucune ligne et ne barre aucun mot
en particulier. Des traits semblables traversent les passages sur les
fig. 79 \emph{b} et 86 aux vues 602 et 605, figures dont les lignes sont
biffées dans les tableaux de la vue 604 ; ils paraissent annuler ces
passages. Le texte est donné tel qu'il est écrit.}
\page{602}

\note{En tête, au milieu : « (63) », barré d'un trait horizontal ; en haut à
droite, à l'encre, d'une écriture plus ronde : 296. Un long trait fin et
oblique descend du haut de la page jusqu'à la ligne « accompagne la
fig. 79 \emph{b} » ; il ne suit aucune ligne et ne barre aucun mot en
particulier. Il couvre, comme ceux des vues 601 et 605, ce qui touche aux
fig. 79 \emph{b} et 86, dont les lignes sont biffées dans les tableaux de
la vue 604, et paraît annuler ce passage ; le texte est donné tel qu'il
est écrit.}

qu'il doit y avoir 6 semblables lignes aulieu de 5 que l'on a trouvées dans le premier cas.
\note{La page précédente s'arrête sur « Mais le terme $\text{sin}\left(\frac{\pi.6y}{A}\right)$ ind », dont la fin
est cachée par l'onglet ; la phrase reprend ici.}

La fig. 86 seroit entièrement conforme à ce résultat si les deux courbes rentrantes qui,
dans la partie superieure de cette figure, remplacent deux des lignes sinueuses,
pouvoient être considérés comme un effet accidentel de quelqu'inégalité dela
plaque. Je serois tenté d'adopter cette opinion, parceque l'existence d'une
ligne sinueuse dans la partie dela plaque qui présente les deux courbes
rentrantes, ne permet guère de supposer qu'il y ait, pour cette moitié, une
relation nouvelle entre $a'$ et $\mathcal{A}$.

Au reste l'application que j'ai cû pouvoir faire de l'intégrale (I)
aux deux figures 79 \emph{b} et 86 n'est que conjecturale ; et il faudroit avoir
recours au calcul pour s'assurer qu'elle satisfait en effet à la description
exacte de ces figures.

En adoptant l'explication que je viens de donner ; on trouvera que le son qui
accompagne la fig. 79 \emph{b}, devra être proportionnel au nombre $m^2=25$ ; et que
le son qui accompagne la fig. 86, devra être proportionnel au nombre $m^2=36$.

On voit que j'ai rendu compte de plus dela moitié des figures observées
par M\textsuperscript{r} Chladni, dans le mouvement des plaques quarrées et que celles
qui n'ont pas été expliqués ne sont ni plus bizares ni plus compliquées
que les autres. Peut-être seroit-il également facile de les expliquer aussi,
soit en mettant de nouvelles valeurs dans les intégrales que j'ai données,
soit en employant alafois plusieurs des mêmes intégrales, comme je le
dirai en parlant des plaques rectangulaires différentes du quarré, soit enfin
a l'aide de quelqu'autre intégrale qui ne s'est pas encore présentée.

Pour completter la comparaison entre la théorie et l'expérience, je vais
m'occuper de la comparaison des sons, et quoique je n'aie encore expliqué
qu'une partie des figures, et que ce ne soit parconséquent que pour cellelà
seulement que je puisse affirmer l'équivalence entre les lignes nodales de
figures variées et les lignes nodales paralleles à un des axes, je suivrai
\page{603}

\note{En tête, au milieu : « (64) », sans rature. Aucun nombre en haut à
droite : c'est le verso du feuillet numéroté 296 à la vue 602. L'écriture
court jusqu'au bord droit du feuillet, où la fin de plusieurs lignes se
perd dans le pli ; les lettres complétées entre crochets d'addition sont
celles que le sens rend certaines.}

à cet égard l'indication de M\textsuperscript{r} Chladni, ensorte que, lorsqu'une figure
n'aura pas été expliquée, je regarderai le son qu'elle produit comme devant
être proportionnel au nombre que la théorie donneroit, s'il s'agissait de
la figure composée de simples lignes droites, à laquelle ce savant
physicien la croit équivalente.

N.o 16 \emph{comparaison des sons}. Si on réfléchit à l'influence que peut avoir
la moindre inégalité d'épaisseur sur les intervalles des sons qui correspon\add{dent}
aux différentes figures que peuvent présenter les plaques élastiques
vibrantes, on ne pourra exiger un accord rigoureux entre les sons
indiqués par la théorie et les sons donnés par l'expérience. Pour
fixer la limite des différences qui devront être négligées entre les
resultats de la théorie et ceux de l'expérience j'observerai que
M\textsuperscript{r} Chladni a déterminé lui même la quantité de l'erreur possib\add{le}
dans les expériences les mieux faites : car on lit p. \uncertain{154} de son livre
« les differences que l'on pourra trouver dans des expériences bien faites
« n'excéderont jamais un demi-ton tout au plus ». Ainsi lorsqu'il
s'agira de comparer entr'eux deux sons donnés par l'expérience, \add{le}
rapport ne pourra jamais différer du rapport entre les nombres auxqu\add{els}
la théorie veut que ces sons soient proportionnels, deplus d'un to\add{n}
encore cette différence pour n'être pas hors des limites dela probabili\add{té}
doit-elle être rare, puisqu'elle suppose que chacune des expériences
comparées, est affectée de la plus grande erreur possible, et que
cette erreur influe, pour chacune d'elles, dans des sens opposés.

Cela posé, je vais commencer par comparer entr'eux les sons donn\add{és}
par les vibrations tournantes, c'est-a-dire par celles qui ne
contiennent pas une période complette, mais seulement deux demies
périodes dans le sens d'un des axes.
\note{« N.o 16 » : son numéro d'article, suivi du titre souligné, donné en
italique ; le 1, court et bas, a l'air d'une virgule après le « o », comme
le 1 de « 10 » et de « 17 » à la vue 604. « p. 154 » : le 5 a la queue descendante qui le fait
ressembler à un 9, et le 4 la forme de « b ». Les guillemets de la citation
sont les siens, répétés en tête de ligne. Après « par l'expérience, » ne
reste qu'un trait au bord du feuillet, où l'article est complété par le
sens.}
\page{604}

\note{En tête, au milieu : « (65) », barré d'un trait horizontal ; en haut à
droite, à l'encre, d'une écriture plus ronde : 297.}

D'après la table donnée p. 152 du traité d'acoustique les fig. qui, parmi celles que j'ai
expliqués, appartiennent à ce cas, sont les suivantes : fig. 63, 66 \emph{a}, \uncertain{69}-70, 74 \emph{a},
79 \emph{a} - \struck{79 \emph{b}}, \struck{86}.

J'ai trouvé n.o 13 que le son de la fig. 63 devoit être proportionnel au nombre 2
\[
\begin{array}{llr}
 & \ldots\ldots\ldots\ldots\ 66\,\underline{a}\ \ldots\ldots\ldots\ldots & 5\\
\text{n.o }15 & \ldots\ldots\ldots\ldots\ 69-70\ \ldots\ldots\ldots\ldots & 10\\
\text{n.o }13 & \ldots\ldots\ldots\ldots\ 74\,\underline{a}\ \ldots\ldots\ldots\ldots & 17\\
\text{n.o }15 & \ldots\ldots\ldots\ldots\ 79\,\underline{a}\ \ldots\ldots\ldots\ldots & 26\\
 & \ldots\ldots\ldots\ldots\ \struck{79\,\underline{b}}\ \ldots\ldots\ldots\ldots & \struck{25}\\
 & \ldots\ldots\ldots\ldots\ \struck{86}\ \ldots\ldots\ldots\ldots & \struck{36}
\end{array}
\]

La difference entre les sons des figures
\[
\begin{array}{lr}
63 \text{ et } 66\,a \ \text{ doit donc être} & \frac{5}{2}\\
63 \text{ et } 69-70\ \ldots\ldots\ldots\ldots & 5\\
63 \text{ et } 74\ \ldots\ldots\ldots\ldots & \frac{17}{2}\\
63 \text{ et } 79\,\underline{a}\ \ldots\ldots\ldots\ldots & 13\\
\struck{63 \text{ et } 79\,\underline{b}}\ \ldots\ldots\ldots\ldots & \struck{\frac{25}{2}}\\
\struck{63 \text{ et } 86}\ \ldots\ldots\ldots\ldots & \struck{18}
\end{array}
\]

La table p. 152 du traité d'acoustique donne : son de la fig.
\[
\begin{array}{ll}
63 \ldots & \text{sol }1\\
66\,a \ldots & \text{si }2\\
69 \ldots & \text{si }3\\
70 \ldots & \text{ut}^{\sharp}\,4\\
74\,\underline{a} \ldots & \text{si}^{\flat}\,4 + \text{grave}\\
79\,a - \struck{79\,b} \ldots & \text{fa}^{\sharp} + \text{grave}\\
\struck{86 \ldots \text{ut }6} &
\end{array}
\]

La difference des sons qui accompagnent les fig. 63 et 66 \emph{a}, est donc un octave
1 ton majeur et un ton mineur, cet intervalle est exprimé, comme on sait, par
$\frac{5}{2}$ c'est-a-dire qu'il est entièrement conforme au resultat de la théorie.

La difference des sons qui accompagnent les fig. 63 et 69, est 2 octaves 1 ton
majeur et un ton mineur. Elle est donc exprimée, comme le veut la théorie par
le nombre 5.

Le son qui accompagne la fig. 70, devroit, d'après la théorie, (v. n.o 15), être
le même que celui qui appartient à la fig. 69. Cependant ces sons different
entr'eux de l'intervalle d'un ton. J'ignore la cause de cette différence,
\note{Les renvois « n.o 13 », « n.o 15 » : le 1, court et bas, se lit
comme une virgule après le « o » ; c'est le même trait que le 1 de « 10 »
et de « 17 » dans le premier tableau. Dans la première liste, le
nombre lu \emph{69}-70 a un 9 sans queue et pourrait se lire 60 ; la table
qui suit porte 69-70. Les lignes relatives aux fig. 79 \emph{b} et 86 sont
barrées d'un trait horizontal dans les trois tableaux et dans la liste,
et sont données ici comme biffées. Les lettres des figures sont soulignées
tantôt oui, tantôt non ; elles sont données comme sur le feuillet. Le
dièse et le bémol sont écrits au-dessus du nom de la note, le chiffre de
l'octave après. « un octave » : ainsi.}
\page{605}

\note{En tête, au milieu : « (66) », sans rature. Aucun nombre en haut à
droite : c'est le verso du feuillet numéroté 297 à la vue 604. L'écriture
court jusqu'au bord droit du feuillet et la fin de plusieurs lignes se perd
dans le pli ; les lettres complétées entre crochets d'addition sont celles
que le sens rend certaines.}

et je ne puis m'empêcher de remarquer combien il est singulier que M\textsuperscript{r} Chla\add{dni}
ait pû se décider à regarder comme équivalentes, deux figures dont il semb\add{le}
que la théorie seul peut indiquer le rapport, puisque, par une circonstan\add{ce}
sans doute accidentelle, l'expérience donne, pour chacune d'elles, des son\add{s}
notablement differens.

La difference entre les sons qui accompagnent les fig. 63 et 74, est de 3
octaves et un peu moins d'un ton et $\frac{1}{2}$ ton, c'est-a-dire qu'elle est
$<$ que $\frac{75}{8}$. En effet $\frac{75}{8}$ est $>$ que $\frac{17}{2}$ donné par la théorie. Les deux rapp\add{orts}
different entr'eux de moins d'un ton et de plus d'un demiton. Ainsi
le son de la fig. 74 semble trop aigu, mais ce cas est pourtant
renfermé dans la limite de l'erreur que l'on peut attribuer aux expérie\add{nces}.

La difference entre les sons qui accompagnent les fig. 63 et 79 \emph{a},
est de 3 octaves et un peu moins de 5 tons et $\frac{1}{2}$ ton, c'est-a-dire $<9$
le nombre donné par la théorie est 13 qui est beaucoup trop petit
car la difference entre les deux rapports est entre 1 ton et 1 ton et
tandis qu'elle devroit être moindre qu'un demiton. Ainsi le so\add{n}
qui accompagne la fig. 79 \emph{a}, semble être trop aigu d'environ 1 ton.

La difference est encore plus grande, lorsqu'il s'agit de la fig. 79 \emph{b} :
la théorie veut que le son qui accompagne cette fig., soit un peu plu\add{s}
grave que le précédent. A la vérité la difference entre les nombre\add{s}
qui appartiennent a ces deux figures savoir $\frac{26}{25}$ est beaucoup moin\add{dre}
qu'un demiton ; et considérée en elle même, elle pourroit être
négligée ; mais elle contribue encore à éloigner la théorie de l'expérie\add{nce}
lorsque l'on prend le son de la fig. 63 pour terme de comparaison.

La difference entre les sons qui accompagnent les figures 63 et 86
est de 4 octaves, d'un ton majeur, d'un ton mineur et d'un $\frac{1}{2}$ ton
majeur : ainsi elle est exprimée par $\frac{64}{3}$. La théorie donne 18, \add{le}
son de la fig. 86 comparé a celui dela fig. 63 est donc plus aigu
\note{« $<9$ » : la ligne s'arrête au bord du feuillet sur ce 9, et la
suite, s'il y en a une, est dans le pli. De même « entre 1 ton et 1 ton
et », dont la fin est perdue. « trop petit » est surchargé. Le 4 de
« 4 octaves » et de « 74 » a sa forme de « b ». Un long trait fin et
oblique part du premier mot de l'alinéa sur la fig. 79 \emph{b} (« La
difference est encore plus grande ») et descend jusqu'au pied de la page,
à travers l'alinéa sur la fig. 86 ; comme ceux des vues 601 et 602, il
couvre ce qui touche aux fig. 79 \emph{b} et 86, dont les lignes sont
biffées dans les tableaux de la vue 604 : il paraît annuler ces passages.
Le texte est donné ici tel qu'il est écrit, sans \emph{biffure}.}
\page{606}

\note{En tête, au milieu : « (67) », barré d'un trait horizontal ; en haut à
droite, à l'encre, d'une écriture plus ronde : 298.}

d'un ton et $\frac{1}{2}$ ton majeur que ne le veut la théorie.

Au reste on se souvient que l'explication que j'ai crû pouvoir donner
des fig. 79 \emph{b} et 86, n'est qu'hypothétique, ensorte que les nombres qui y
correspondent, ne sont que probables.
\note{Le trait fin et oblique de la vue 605 reprend ici et traverse les
quatre premières lignes, jusqu'à « que probables » ; il couvre la fin de
l'alinéa sur la fig. 86 et la remarque sur les fig. 79 \emph{b} et 86.}

Je vais m'occuper à present de comparer les sons qui appartiennent aux
cas où les figures sont composées d'une ou de plusieurs périodes complettes
dans le sens de chacun des deux axes.

Je commencerai par celles de ces fig. dont j'ai rendu compte dans les n.os précédens,
et je prendrai pour terme de comparaison la fig. 68 \emph{a} qui est la plus
simple de cette classe.

J'ai trouvé n.o 13 que le son de la fig. 68 \emph{a} devoit être proportionnel au nombre 8
\[
\begin{array}{llr}
 & \ldots\ldots\ldots\ldots\ 71\ \ldots\ldots\ldots\ldots & 13\\
 & \ldots\ldots\ldots\ldots\ 75\ \ldots\ldots\ldots\ldots & 18\\
 & \ldots\ldots\ldots\ldots\ 80\,\underline{a}\ \ldots\ldots\ldots\ldots & 25\\
 & \ldots\ldots\ldots\ldots\ 81\,\underline{a}\ \ldots\ldots\ldots\ldots & 29\\
 & \ldots\ldots\ldots\ldots\ 82\ \ldots\ldots\ldots\ldots & 32\\
\text{n.o }15 & \ldots\ldots\ldots\ldots\ 87\,\underline{a},\ 87\,\underline{b},\ 88\,\underline{a} \text{ et } 88\,\underline{b}\ \ldots\ldots & 40\\
 & \ldots\ldots\ldots\ldots\ \struck{79\,\underline{a}}\ \ldots\ldots\ldots\ldots & \struck{26}
\end{array}
\]

La difference entre les sons des fig.
\[
\begin{array}{lr}
68 \text{ et } 71 \ \text{ doit donc être} & \frac{13}{8}\\
68 \text{ et } 75\ \ldots\ldots\ldots\ldots & \frac{9}{4}\\
68 \text{ et } 80\,\underline{a}\ \ldots\ldots\ldots\ldots & \frac{25}{8}\\
68 \text{ et } 81\,\underline{a}\ \ldots\ldots\ldots\ldots & \frac{29}{8}\\
68 \text{ et } 82\ \ldots\ldots\ldots\ldots & 4\\
68 \text{ et } 87-88\ \ldots\ldots\ldots\ldots & 5\\
\struck{68 \text{ et } 79\,\underline{a}}\ \ldots\ldots\ldots\ldots & \frac{13}{4}
\end{array}
\]

La table p. 152 du traité d'acoustique : son de la fig.
\[
\begin{array}{ll}
68 \ldots & \text{si}^{\flat}\,3 - \text{aigu}\\
71 \ldots & \text{fa}^{\sharp}\,4\\
75 \ldots & \text{ut }5\\
80 \ldots & \text{fa}^{\sharp}\,5\\
81 \ldots & \text{sol}^{\sharp}\,5 + \text{aigu}\\
82 \ldots & \text{si}^{\flat}\,5\\
87-88 \ldots & \text{ut}^{\sharp}\,6 + \text{aigu}\quad \text{re }6 + \text{grave}\\
\struck{79 \ldots \text{fa}^{\sharp}\,5 + \text{grave}} &
\end{array}
\]
\note{La dernière ligne du premier tableau, la ligne « 68 et 79 \emph{a} » du deuxième (dont la
fraction $\frac{13}{4}$ n'est pas barrée) et la dernière ligne du troisième,
relatives à la fig. 79 \emph{a}, sont barrées d'un trait horizontal. Le 1
de « n.o 13 », « n.o 15 » est, comme à la vue 604, un trait court et bas.
Devant « aigu », à la première ligne de la table, un tiret et non le « + »
des autres lignes. « re 6 » : sans accent.}
\page{607}

\note{En tête, au milieu : « (68) », sans rature. Aucun nombre en haut à
droite : c'est le verso du feuillet numéroté 298 à la vue 606. Plusieurs
fins de ligne se perdent au bord droit du feuillet ; les lettres complétées
entre crochets d'addition sont celles que le sens rend certaines.}

La difference entre les sons qui accompagnent les fig. 68 et 71, est donc un
peu plus de 4 tons, c'est a dire $>$ que $\frac{8}{5}$ la théorie donne $\frac{13}{8}$ qui est en effet
un peu $>$ que $\frac{8}{5}$.

La difference entre les sons qui accompagnent les fig. 68 et 75, est un octave, et u\add{n}
peu plus d'un ton c'est-a-dire $>\frac{9}{4}$ qui est précisément le nombre donné par l\add{a}
théorie.

La difference entre les sons qui accompagnent les fig. 68 et 80, est d'un octa\add{ve}
et un peu plus de 4 tons, c'est-a-dire $>$ que $\frac{16}{5}$. La théorie donne
au contraire $\frac{25}{8}$ qui est $<$ que $\frac{16}{5}$ : mais la difference entre les deux
rapports, est bien audessous de la limite des erreurs puisqu'elle
est juste un \emph{comma}.

La difference entre les sons qui accompagnent les fig. 68 et 81 est
notablement plus grande qu'un octave et 5 tons c'est-a-dire $>$ que
et en effet le rapport donné par la théorie savoir $\frac{29}{8}$ surpasse
$\frac{32}{9}$ de plus d'un \emph{comma}.

La difference entre les sons qui accompagnent les fig. 68 et 82 est un peu $>$
2 octaves, c'est-a-dire $>$ que 4 qui est le nombre donné par la théorie.

La difference entre les sons qui accompagnent les fig. 68 et 87-88, le
different entr'eux de moins d'un demi ton, est, en adoptant le plu\add{s}
favorable a la théorie, de deux octaves et de deux tons : elle est donc
exprimée par le nombre 5 donné en effet par la théorie.

La difference entre les sons qui accompagnent les fig. 68 et 79,
est un octave 3 tons majeurs et 1 ton mineur : elle est donc
exprimée par $\frac{16}{5}$. La théorie donne $\frac{13}{4}$ qui ne differe pas sensiblem\add{ent}
de ce rapport, puisque $\frac{65}{64}$ est $<$ que le \emph{comma} $\frac{128}{125}$.

Les cas suivans sont ceux pour lesquels M\textsuperscript{r} Chladni n'a pas do\add{nné}
de fig. et aussi ceux dont les fig. n'ont pas été expliquées dans les
\note{« c'est-a-dire $>$ que » (fig. 68 et 81) : la ligne s'arrête au bord du
feuillet après « que », et la valeur, s'il y en a une, est perdue dans le
pli. « le » : au bord du feuillet, la fin du mot est perdue dans le pli (le sens demande « lesquels »).
« est un peu » (fig. 68 et 82) est écrit plus petit, serré au bout de la
ligne. « de deux octaves » : « deux » est écrit sur un autre mot, peut-être
« d'un ». Un long trait fin et oblique traverse tout l'alinéa sur les
fig. 68 et 79, du premier mot à « comma » ; il répond à la ligne barrée
« 68 et 79 \emph{a} » de la vue 606 et paraît annuler l'alinéa, donné ici
tel qu'il est écrit. Avant « comma », « le » est écrit sur « les ».}
\page{608}

\note{En tête, au milieu : « (69) », barré d'un trait horizontal ; en haut à
droite, à l'encre, d'une écriture plus ronde : 299.}

n.os précédans. Pour les désigner, j'emploirai la notation imaginée par cet auteur,
suivant laquelle $m|n$ signifira qu'il devra y avoir \emph{m} lignes droites nodales
parallelles à un des axes et \emph{n} pareilles lignes parallelles à l'autre axe.

Je continurai de comparer les sons des autres cas à celui qui accompagne
la figure 68 \emph{a} c'est adire a celui du cas $2|2$. Cela posé, la théorie montre
que pour $2|2$ le son doit être proportionnel à 8
\[
\begin{array}{lr}
4|2\ \ldots\ldots\ldots\ldots\ldots\ldots\ldots\ldots & 20\\
5|3\ \ldots\ldots\ldots\ldots\ldots\ldots\ldots\ldots & 34\\
6|3\ \ldots\ldots\ldots\ldots\ldots\ldots\ldots\ldots & 45\\
5|4\ \ldots\ldots\ldots\ldots\ldots\ldots\ldots\ldots & 41\\
6|4\ \ldots\ldots\ldots\ldots\ldots\ldots\ldots\ldots & 52\\
5|5\ \ldots\ldots\ldots\ldots\ldots\ldots\ldots\ldots & 50\\
6|5\ \ldots\ldots\ldots\ldots\ldots\ldots\ldots\ldots & 61
\end{array}
\]

La difference entre
\[
\begin{array}{lr}
2|2 \text{ et } 4|2 \ \text{ doit donc être} & \frac{5}{2}\\
2|2 \text{ et } 5|3\ \ldots\ldots\ldots\ldots & \frac{17}{4}\\
2|2 \text{ et } 6|3\ \ldots\ldots\ldots\ldots & \frac{45}{8}\\
2|2 \text{ et } 5|4\ \ldots\ldots\ldots\ldots & \frac{41}{8}\\
2|2 \text{ et } 6|4\ \ldots\ldots\ldots\ldots & \frac{13}{2}\\
2|2 \text{ et } 5|5\ \ldots\ldots\ldots\ldots & \frac{25}{4}\\
2|2 \text{ et } 6|5\ \ldots\ldots\ldots\ldots & \frac{61}{8}
\end{array}
\]

La table p. 152 donne : son de
\[
\begin{array}{ll}
2|2 \ldots & \text{si}^{\flat}\,3 + \text{grave}\\
4|2 \ldots & \text{ut}^{\sharp}\,5 \text{ et re }5\\
5|3 \ldots & \text{si }5 + \text{grave et ut }6 + \text{grave}\\
6|3 \ldots & \text{mi }6\\
5|4 \ldots & \text{re}^{\sharp}\,6\\
6|4 \ldots & \text{sol}.\,6 + \text{aigu}\\
5|5 \ldots & \text{fa}^{\sharp}\,6 + \text{aigu}\\
6|5 \ldots & \text{si}^{\flat}_{6} + \text{grave}
\end{array}
\]

On voit que, les deux cas qui appartiennent à $4|2$ different entr'eux
d'un demiton si on prend re 5 + grave, cequi rapproche d'ut$^{\sharp}$ 5, on trouvera
que la difference entre les sons de $2|2$ et de $4|2$, sera d'un octave
et 2 tons, c'est-a-dire qu'elle sera exprimée comme le veut la théorie par $\frac{5}{2}$.
\note{« n.os précédans » achève la phrase de la vue 607 (« dont les fig.
n'ont pas été expliquées dans les »). Dans la notation de Chladni, $m|n$ est
écrit avec un trait vertical entre les deux nombres ; \emph{m} et \emph{n}
sont soulignés par elle. À la dernière ligne de la table, le 6 est écrit
sous le bémol de « si ». « qui » (avant « appartiennent ») est surchargé.}
\page{609}

\note{En tête, au milieu : « (70) », sans rature. Aucun nombre en haut à
droite : c'est le verso du feuillet numéroté 299 à la vue 608.}

Pour $5|3$, les deux cas different entr'eux d'un ton environ. Leplus favorable à
la théorie, savoir si 5 + grave, donne dans la comparaison à $2|2$, une difference
de deux octaves et $\frac{1}{2}$ ton. Elle est donc exprimée par $\frac{25}{6}$ ou par $\frac{64}{15}$ ;
et $\frac{17}{4}$ qui est le rapport indiqué par la théorie, est en effet $>$ que $\frac{25}{6}$
et $<$ que $\frac{64}{15}$.

La difference entre les sons qui appartiennent à $2|2$ et a $6|3$ est d'un
peu plus de deux octaves et 3 tons, c'est-a-dire $>$ que $\frac{50}{9}$ ; et en effet
le rapport donné par la théorie savoir $\frac{45}{8}$ est de fort peu $>$ que $\frac{50}{9}$.

La difference entre les sons qui appartiennent à $2|2$ et à $5|4$, est
un peu plus de 2 octaves 2 tons et $\frac{1}{2}$ ton, c'est-a-dire $>$ que $\frac{128}{25}$ ;
et en effet le rapport donné par la théorie, savoir $\frac{41}{8}$ est $>$ que
$\frac{128}{25}$ ; mais d'une quantité absolument inappréciable a l'oreille.

La difference entre les sons qui appartiennent à $2|2$ et à $6|4$,
est plus de 2 octaves 4 tons et $\frac{1}{2}$ ton, c'est-a-dire $>$ que $\frac{20}{3}$.
Aucontraire le rapport donné par la théorie savoir $\frac{13}{2}$ est $<$ que
$\frac{20}{3}$ : mais la difference entre ces deux rapports, est fort petite,
puisqu'elle surpasse le \emph{comma} d'une quantité inapréciable a l'oreille.

La difference entre les sons qui appartiennent à $2|2$ et à $5|5$, est un
peu plus grande que 2 octaves et 4 tons, c'est-a-dire $>$ que $\frac{32}{5}$.
Aucontraire le rapport donné par la théorie, savoir $\frac{25}{4}$, est $<$ que
$\frac{32}{5}$ : mais la difference entre ces deux nombres est bien petite,
puisqu'elle se trouve juste d'un \emph{comma}.

La difference entre les sons qui appartiennent à $2|2$ et à $6|5$,
est de 3 octaves ; elle est donc exprimée par le nombre 8 qui differe
de moins d'un demi ton majeur du nombre $\frac{61}{8}$ donné par la théorie.

Il est impossible d'espérer plus d'accord entre la théorie et
l'expérience, que n'en présente la comparaison des sons qui accompagne\add{nt}
\note{Au premier $\frac{128}{25}$, un trait précède le 25 du dénominateur, qu'on
pourrait prendre pour un 3 ; le second est net. Le premier
« inappréciable » est surchargé ; le second est écrit « inapréciable ».}
\page{610}

\note{En tête, au milieu : « (71) », barré d'un trait horizontal ; en haut à
droite, à l'encre, d'une écriture plus ronde : 300.}

celles des figures que j'ai expliqués, dans lesquelles il y a une ou plusieurs
périodes complettes, tant dans le sens des \emph{x} que dans celui des \emph{y}.

Le même accord se trouve aussi dans la comparaison des sons que, d'après
l'autorité de M\textsuperscript{r} Chladni, j'ai attribués aux cas des périodes complettes.
En effet, en prenant dans cette dernière classe, qui présente d'assez grandes
differences de sons rapportés aux mêmes cas, ceux de ces \struck{sons} \add{sons} qui s'accordent
le mieux avec la théorie, j'ai trouvé que les rapports indiqués par
la théorie, differoient toujours de moins d'un demi ton de ceux donnés
par l'expérience. On voit que cette difference est bien audessous de
celles dont les erreurs possibles pourroient rendre compte ; et que M\textsuperscript{r}
Chladni a plusieurs fois regardé comme équivalentes, des figures qui
donnoient des sons differens entr'eux d'un demi ton, d'un ton, ou même
d'une tierce (v. la table p. 152 du traité d'acoustique).

Cependant il se présente ici une difficulté : lorsqu'au lieu de comparer
entr'eux les sons donnés par les figures qui renferment des périodes
complettes dans le sens des \emph{x} et dans celui des \emph{y}, on prend pour terme
de comparaison les sons qui accompagnent les figures qui appartiennent
à l'intégrale (C) et qui, tant d'après l'opinion de M\textsuperscript{r} Chladni
que d'après la théorie que j'ai cherché a établir, doivent être les
mêmes que les sons donnés pour la simple lame dont les extrémités sont
libres en vertu des conditions $\frac{ddz}{dx^2}=0$ et $\frac{d^3z}{dx^3}=0$, on trouve que
le son de toutes les figures périodiques qui renferment des périodes
complettes, est plus grave d'environ 1 ton que la théorie ne l'indique.
C'est ainsi, par exemple, que le son de $0|4$ fig. 72-73 est sol$^{\sharp}$ 3
et que celui de $4|3$ fig. 80 est fa$^{\sharp}$ 5, tandis qu'il ne devroit y avoir que
deux octaves entre ces deux sons, puisque, d'après la théorie d'Euler,
le premier est proportionnel au nombre $\frac{25}{4}$, et que d'après la théorie
des surfaces périodiques, le second est proportionnel au nombre 25.
\note{Les lettres \emph{x} et \emph{y} sont soulignées par elle. « sons » est
biffé et récrit au-dessus de la ligne. « de toutes » : le mot est
surchargé.
« Euler » est écrit avec son E en forme de L.}
\page{613}

\note{En tête, au milieu : « (72) », sans rature. Aucun nombre en haut à
droite : c'est le verso du feuillet numéroté 300 à la vue 610.}

Mais le son des figures qui contiennent ou sont censées contenir une ou
plusieurs périodes complettes, ne sont pas seulement plus graves d'un ton
environ que ne semble l'exiger la comparaison aux sons donné par la
simple lame ; cette difference doit encore être considérée comme un
\emph{minimum} dont se rapprochent d'autant plus les sons produis par les
vibrations \emph{tournantes} que les figures qui leur appartiennent contiennent
un plus grand nombre de lignes nodales. C'est ainsi que d'une part
les sons qui accompagnent les fig. 63 et 68, bien qu'ils ne devroient
differer que de deux octaves puisqu'ils sont entr'eux, d'après la
théorie $::2:8$ different pourtant de 2 octaves 1 ton et $\frac{1}{2}$ ton : desorte que
le son de la fig. 63 est plus grave de 2 tons et $\frac{1}{2}$ ton que ne le
veut la comparaison aux sons de la simple lame ; et que les sons
qui accompagnent les fig. 66 \emph{a} et 69, présentent les mêmes differences,
tandis que, d'autre part, le son de la fig. 74 est moins éloigné du
rapport donné par la théorie ; et que le son dela fig. 79 \struck{\emph{b} et 86}
se rapproche encore \struck{plus} \add{d'avantage du rapport} \struck{dela nombres} \add{donné par} \struck{même} théorie puisqu'il n'en differe
\add{plus que de l'intervalle d'un ton et $\frac{1}{2}$ ton.} que j'ai crû pouvoir leur
\struck{assigner en leur appliquant l'intégrale (I)}, car le son dela fig. 8
qui \struck{est ut 6 comparé à celui de la fig. 75 qui est ut 5
(v. la table p. 152) est exactement cequ'il doit être pour la comparaison
aux cas des périodes complettes, puisque le nombre que j'ai assigné
au premier est 36 et que la théorie donne 18 pour le second.}

Des differences du même genre se font remarquer dans la compar\add{aison}
des sons des figures qui se forment sur les plaques rectangulaires
differentes du quarré desorte que les figures qui sur ces diverses
plaques présentent une ou plusieurs périodes complettes donnent
toujours, au moins en prenant un résultat moyen, des sons un
peu plus graves qu'il ne le faudroit pour établir l'accord dela
théorie et de l'expériencedans la comparaison au son des figures
\note{« le son », « se rapproche » : elle a écrit « les sons », « se
rapprochent », puis annulé le pluriel par son petit « x » sur l's final
(« lex sonx ») et en barrant la fin du verbe ; le singulier est donné.
« plus que de l'intervalle d'un ton et $\frac{1}{2}$ ton. » est écrit dans
l'interligne, sous la ligne corrigée. La fin de la phrase primitive, de
« que j'ai crû pouvoir leur » à « pour le second », est à annuler avec
elle : la plupart des mots en sont barrés d'un trait, mais « que j'ai crû
pouvoir leur », « car le son dela fig. 8 » et « qui » ne portent pas le
trait ; ils sont donnés ici tels qu'ils sont. Le numéro de figure après
« fig. 8 » se perd au bord du feuillet (86, par la suite de la phrase).
« assigner », « assigné » : son double s ressemble à « tt ». « produis
par » : ainsi. « l'expériencedans » : en un mot. \emph{minimum} et
\emph{tournantes} sont soulignés par elle.}
\page{614}

\note{Autre pièce. Feuillet d'un autre papier, sans pagination, d'une
écriture plus rapide, corrigée dans les interlignes : un brouillon sur
l'équation des surfaces élastiques, et non une page du mémoire paginé qui
s'arrête à la vue 613. En haut à droite, à l'encre, de l'écriture ronde de
la série : 301. La première ligne commence au milieu d'une phrase ; son
début n'est pas sur les vues voisines de ce lot.}

\struck{\ill{}} differentielles des deux ordonnées rectangulaires \emph{x} et \emph{y} les differentielles des deux
\uncertain{courbures} \struck{pareillement rectangulaires} \add{qui mesurent les courbures principales} dela surface desorte que \struck{la} \add{les} valeur\struck{s} de \emph{z}
\struck{n'exprimeroit encore} \add{roient}, comme dans le cas du plan, la quantité dont chaque
point de cette surface s'éleveroit pendant le mouvement, audessus dela situation
naturelle.

L'équation que je crois appartenir au mouvement de vibration des surfaces
élastique est comme on \uncertain{valevoir} du quatrième dégré, elle ne peut avoir
d'intégrale generale en termes finis mais elle peut être abaissée au second
dégré et dans le cas des mouvemens comparables a ceux du pendule
l'une des équations auxquelles elle se réduit n'est autre chose que
l'équation des surfaces tendues sous la même condition. \struck{il resulte}
\uncertain{delà} ainsi ladifficulté de trouver des intégrales particulières convenables
aux differens cas du mouvement des surfaces elastiques se trouve réduite
a obtenir celles des \struck{mêmes cas \ill{}} surfaces tendues car il est evident
qu'en joignant les intégrales de chacune des \struck{surf} équations du
second degré on aura une intégrale de l'équation du quatrième
degré \struck{du degré de généralité que comporte} et que cette intégrale
aura toute la généralité que comporte l'intégrale de l'équation aux
surfaces tendues dont elle est composée, on voit \struck{aussi} \add{encore} que cette
dernière intégrale peut aussi convenir seule a l'équation du quatrième
degré desorte que \struck{toute figure que peut prendre} \add{quelque soit la figure que puisse prendre} une surface
tendue \struck{sous les mêmes} conditions \add{renfermée dans des limites données la même figure peut egalement}, appartenir, a une surface
élastique renfermée dans les mêmes limites.
\note{Les corrections se lisent ainsi : au-dessus de « pareillement
rectangulaires », barré, « qui mesurent les courbures principales » ; au-dessus
de « la valeur », « les » écrit sur « la », et l's de « valeurs » annulé par
son petit « x » ; au-dessus de « n'exprimeroit », barré avec « encore »,
la terminaison « roient », sans que le verbe corrigé soit récrit. Le mot lu
\uncertain{courbures}, en tête de la deuxième ligne, a sa fin surchargée ;
le premier mot de la page est barré et n'a pas été lu. « comme on
\uncertain{valevoir} » : « va le voir » en un mot, lu avec doute. Le
début de « delà » est lu avec doute, au commencement de la ligne. Le mot
barré après « des mêmes cas » n'a pas été lu. « l'équation » (avant « aux surfaces
tendues ») est écrit sur un autre mot. Les ajouts « encore », « quelque soit la figure
que puisse prendre » et « renfermée dans des limites données la même
figure peut egalement » sont dans l'interligne, au-dessus des mots barrés ;
« sous les mêmes » est barré, « conditions » ne l'est pas, mais la phrase
corrigée l'exclut. Le bas du feuillet est blanc.}
\page{616}

\note{Feuillet du même papier et de la même main rapide que celui de la vue
614, sans pagination ; en haut à droite, à l'encre : 302. Il commence un
nouvel alinéa et ne continue pas visiblement la phrase de la vue 614, dont
le feuillet s'achève sur un point.}

Les geometres ne se sont pas encore beaucoup occupés dela détermination
de ces intégrales \struck{\uncertain{ni en general}} \add{et encore moins} de l'examen des differentes figures
nodales que peuvent présenter \add{la vibration} des surfaces tendues \struck{mises en mouvement
en vibration vibrantes}. M\textsuperscript{r} Biot est je crois le seul qui ait démontré
l'éxistence des lignes nodales parallelles a chacun des côtés de\struck{s} \add{telles} \struck{mêmes}
surfaces rectangulaires \struck{mais comme \uncertain{ce savant} géometre ne s'est occupé
qu'en passant de cet objet il a borné ses recherches a ce cas simple}
je montrerai que cette propriété est une conséquence nécessaire de la
périodicité de ces surfaces pour lesquelles il existe entre les trois coordonnées
rectangulaires une équation de la forme de celle aux cordes vibrantes
et je donnerai plusieurs intégrales particulières dont les résultats
seront ensuite comparés a l'expérience \struck{\uncertain{d'après les le plus}} ainsi
quoique l'on ait encore fait aucun essais sur les figures nodales
que présenteroit les differens cas du mouvement vibratoire des
surfaces tendues on pourra connoître avec certitude quelles seroient
ces figures et aquel son elles correspondroit par la comparaison
des formules aux cas correspondans des surfaces elastiques

Une singularité remarquable est que ces cas ou les surfaces élastiques
prennent dans leur mouvemens vibratoires des figures qui \struck{correspondroient}
pourroient appartenir egalement a des surfaces tendues sont les seuls
dans lesquels il soit permis de considérer les \emph{nappes vibrantes}
comme analogues aux \emph{ondes} des cordes de Sauveur et les \emph{courbes
d'équilibre} ou de \emph{repos} comme correspondans aux \emph{nœuds} ou
\emph{points de repos} des mêmes cordes. dans tout autre cas cette analogie
si naturelle a présumer et indiquée par les termes mêmes du
\note{« ni en general » : lu avec doute sous la rature ; « et encore moins »
est écrit au-dessus, « moins » lui-même barré d'un trait léger. « la
vibration » est ajouté au-dessus de « présenter », sans que « des » soit
corrigé. « de\struck{s} \add{telles} \struck{mêmes} » : « telles » est écrit
au-dessus de « mêmes », barré ; l's de « des » est surchargé. Dans le
passage barré « mais comme ce savant géometre ne s'est occupé qu'en
passant de cet objet… ce cas simple », « ce savant » est lu avec doute et
les derniers mots se perdent sous la rature. « d'après les le plus » : lu
sous la rature, avec doute. « aucun essais », « présenteroit »,
« correspondroit » : ainsi. Les mots soulignés par elle sont donnés en
italique ; « nœuds » est écrit « naudz », avec l'œ ouvert et son petit
« x » final. Le bas du feuillet est blanc et la phrase reste suspendue sur
« les termes mêmes du ».}
\page{618}

\note{Feuillet du même papier et de la même main que ceux des vues 614 et
616, sans pagination ; en haut à droite, à l'encre : 303. La première ligne
achève la phrase de la vue 616 (« indiquée par les termes mêmes du »). Le
texte n'occupe que le haut du feuillet ; le reste est blanc.}

programme \add{publié par la classe} ne se verifie pas parceque comme je le dirai dans
la suite de ce mémoire, il n'est plus vrai alors que les lignes
de repos puissent etre considérées comme limites des parties vibrantes
qu'elles separent.

Après avoir exposé la manière dont je suis parvenu a l'équation
differentielle des surfaces elastiques je donnerai differentes intégrales particulières
dont je comparerai ensuite les resultats al'experience je m'estimerai
heureux si la classe juge que j'ai reussi a ebaucher une théorie
que des
\note{« publié par la classe » est écrit au-dessus de « programme », sans
signe d'insertion. « parvenu », « heureux » : au masculin sur le feuillet.
La phrase s'arrête sur « que des », au milieu de la ligne.}
\page{620}

\note{Feuillet du même papier et de la même main que ceux des vues 614 à
618, sans pagination au milieu ; en haut à droite, à l'encre : 304, et à
sa droite un petit chiffre noirci, peut-être 1, barré d'un trait oblique et
souligné d'un trait ondulé : une numérotation qu'elle a annulée. Le texte
ne continue pas la phrase de la vue 618 (« une théorie que des »), qui
reste suspendue. Le bas du feuillet est blanc.}

Après les cas dont je viens de parler le plus simple auquel on puisse appliquer
l'équation (e) est celui d'une surface cylindrique a base circulaire
En prenant $R$ pour le rayon de cette base l'équation (e) deviendra
\[
\kappa b^{2}_{\ill{}}\Bigl\{\frac{d^4z''}{dx^4}+2\frac{d^4z''}{dx^2dy^2}+\frac{d^4z''}{dy^4}\ \ill{}\ \frac{1}{R^2}\Bigl\{\frac{d^2z''}{dx^2}+\frac{d^2z''}{dy^2}\Bigr\}\Bigr\}=z''
\]
Si aulieu del'équation dela surface entière on veut avoir à présent celle
d'une portion quelconque $A$ dela même surface \struck{et qu'on nomme} \add{soit} $A$
l'arc de cette portion et $\underline{a}$ l'arc correspondant \add{a $A$ dans} le cercle dont le rayon
est pris pour unité on aura la proportion
\[
1:R::a:A \quad\text{on en tire}\quad R=\frac{A}{a}\,.\ \frac{1}{R^2}=\frac{a^2}{A^2}.
\]
Le genre de surface auquelles appartiendroit l'équation
\[
\kappa b^{2}_{\ill{}}\Bigl\{\frac{d^4z''}{dx^4}+2\frac{d^4z''}{dx^2dy^2}+\frac{d^4z''}{dy^4}+\frac{a^2}{A^2}\Bigl\{\frac{d^2z''}{dx^2}+\frac{d^2z''}{dy^2}\Bigr\}\Bigr\}=z''
\]
n'a encore, je crois, été l'objet d'aucune experience.

\struck{Le tems m'a manqué pour effectuer toutes celles que j'avois projettés.
Le desir d'appliquer la théorie que renferme d'ajouter la}
ainsi je me \struck{suis} trouvé \add{privée ici} de tout secours \add{étranger} \struck{pour établir sur des faits
avant connu la théorie qu'elle renferme.}

Je me suis \struck{avisé} \uncertain{résolue} beaucoup trop tard a reprendre le \struck{travail} \add{sujet}
qui fait l'objet du présent mémoire le tems me manque
alafois \add{et} pour \struck{faire} exécuter la totalité des experiences que j'aurois
desiré faire et pour m'appliquer a la recherche des
intégrales particulières qui devroient rendre raison des
\struck{leur} resultats de ces experiences

Quelqu'imparfait que soit cette partie de mon travail je
vais exposer ici \struck{l'état de mes connaiss} les tentatives que j'ai
faits et le succès que j'ai obtenu. Je commencerai
par indiquer les resultats generaux des experiences
je chercherai \struck{ensuite} dans la théorie la manière d'en
rendre compte. \struck{de ces} des singularités qu'ils présente
\note{Les deux équations sont écrites au centre de la ligne ; les accolades
sont les siennes, la fermante tracée comme une virgule doublée.
Devant le coefficient, un $\kappa$ (ou un $k$
à boucle), puis $b$ avec un 2 en exposant et, sous ce 2, un second petit
signe qui n'a pas été lu. Dans la première équation, le signe avant
$\frac{1}{R^2}$ est noirci d'une tache ; la seconde porte $+$ à la même
place. Dans $\frac{1}{R^2}=\frac{a^2}{A^2}$ le 1 est écrit sur un autre
chiffre et le $A$ du dénominateur est surchargé ; au-dessus de $\frac{a^2}{A^2}$
de la seconde équation, le $a$ est lui aussi repassé. « soit » est écrit
au-dessus de « nomme », barré avec « et qu'on » ; « a $A$ dans » est
écrit au-dessus de « le cercle ». « auquelles appartiendroit » : la
terminaison « droit » est écrite sur une autre. Les deux lignes « Le tems
m'a manqué… » et « Le desir d'appliquer… » sont barrées ; dans la ligne
suivante, « privée ici » est écrit au-dessus de « suis », barré, et de
« trouvé », et « étranger » au-dessus de « pour établir », barré avec la
fin de la phrase : la phrase corrigée se lit « ainsi je me trouvé privée
ici de tout secours étranger ». « \uncertain{résolue} » : lu avec doute,
le mot est serré après « avisé » barré. « privée », « résolue » : au
féminin, alors que « parvenu » et « heureux » (vue 618) sont au masculin.
« l'état de mes connaiss » : lu sous des ratures en croix. Dans « des experiences », la fin de « des » est surchargée. La page s'arrête sur « qu'ils
présente », au milieu de la ligne.}

\end{document}
